


\chapter{Related Work}


Two research communities built the foundation of this thesis, and we bring them together throughout the following chapters: algorithmic fairness (documenting demographic harm) and the science of evaluations (challenging the validity of measurements). In the following sections, we outline the findings in the respective fields and position our work at the intersection: investigating current evaluation methods and building new ones that allow us to systematically measure safety risks modulated by user demographics.



\section{Normative Framings of Safety}


From gender and racial biases in vision models \cite{Buolamwini2018GenderSI} to similar biases in today's language models \cite{parrish-etal-2022-bbq}, the field of AI Ethics has long documented that AI models exhibit demographic biases. AI Safety, as the field focussed on large-scale and individual harm from LLMs, developed largely in parallel. However, the normative framing of what \textit{ethics} versus \textit{safety} means for AI is fuzzy itself, and the distinction collapses once we start asking whether safeguards protect all users equally. As Lazar and Nelson emphasise \cite{lazar2023ai-a8d}, \textit{AI Safety} cannot be a purely neutral technical target but inherently contains normative choices about whose values and harms count. \textit{Safety} requires a holistically sociotechnical approach, if the goal is to protect all humans from AI risks.
 Furthermore, the safety community's frequent insistence on prioritising existential risks from AI may miss the widespread and potentially equally catastrophic harm from everyday risks accumulating through gradual disempowerment \cite{kulveit2025gradualdisempowermentsystemicexistential}. Everyday safety for diverse user populations may hence be a defensible core of AI safety rather than a distraction from safety concerned with global catastrophic risk \cite{gyevnar2025ai-cae}.
 More concretely, AI alignment (the problem of ensuring AIs follow, endorse, and act on a desired set of values) must contend with the plurality of human values, which directly necessitates a pluralistic approach to AI development and evaluation \cite{sorensen2024roadmap}. In Chapter \ref{chap:evaleval}, we will hence investigate to which degree current safety evaluations follow a pluralistic approach and evaluate safety concerns against diverse user populations.


\section{Algorithmic Fairness and Demographic Biases in LLMs}


\subsection{Harm Taxonomies and Typologies}

Blodgett et al.  \cite{blodgett-etal-2020-language} critique much of existing work regarding fairness in LLMs for failing to state what is harmful, in what way, to whom and why. We aim to clarify the \textit{what} and \textit{in what way} based on the literature and the \textit{to whom} and \textit{why} through our own investigations for safety evaluations:

The MIT AI Risk Repository \cite{slattery2024airisk}, an aggregated meta taxonomy of more than 65 AI risk frameworks, lists demographic discrimination as a first class risk domain among its seven main risk categories (Discrimination \& Toxicity, Privacy \& Security, Misinformation, Malicious Actors, Human-Computer Interaction, Socioeconomic \& Environmental, AI System Safety, Failures \& Limitations). More explicitly, Weidinger et al. \cite{weidinger2022taxonomy} state "Lower performance for some languages and social groups" inside their discrimination category.
While these taxonomies establish the \textit{what is harmful}, Shelby et al. \cite{shelby2023sociotechnical} address the \textit{in what way} with a typology of how demographically-mediated performance and safety gradients matter. Representational harm occurs when models exhibit biases that reinforce users' stereotypical beliefs or systematically misrepresent certain population groups to their disadvantage. Beyond representational harm, Shelby et al. argue for \textit{quality-of-service} harm in which certain user groups experience degraded quality when using AI systems. Similarly, Solaiman et al. \cite{solaiman2023social} explicitly call for evaluations of "disparate performance" as a key evaluation target for social impact of AI systems. Chehbouni et al. \cite{chehbouni-etal-2024-representational} provide an empirical antecedent for this latter phenomenon showing that Llama 2 Chat's safety safeguards have disparate behaviours dependent on demographics like religion or ethnicity when it comes to overrefusal.
This thesis aims to build on Chehbouni et al. by applying and investigating the quality-of-service disparities for safety specifically.

\subsection{Precedents of Demographics-driven Quality-of-Service Disparities}

In the vein of aligning models to human values, one might hope that alignment techniques reduce biases as values of fairness and equality get encoded. However, as shown by Hofmann et al. \cite{hofmann2024ai-e4f}, covert, dialect-conditioned prejudices, in this case against African American English, survive alignment training through Reinforcement Learning from Human Feedback and even worsen in larger models. We will discuss linguistic diversity in Chapter \ref{chap:translations} in detail. Besides dialect-driven disparities, Salinas et al. \cite{salinas2024name} observe name-driven shifts in LLM-generated general advice and outcomes in audits that vary race- and gender-coded names. An et al. replicate the finding for hiring-style decisions \cite{an-etal-2024-large}. Our concern about user-identity driven changes in the quality of the response, which we will address in Chapter \ref{chap:contextualising}, is further supported by Kantharuban et al. \cite{kantharuban-etal-2025-stereotype}, who demonstrate that LLM-generated recommendations reflect both what the user wants and who the user is, but that models fail to transparently disclose that the latter influenced the recommendation. Likewise, Eloundou et al.'s large-scale "first-person fairness" audits operationalise fairness towards the user for general everyday use cases \cite{eloundou2025firstperson}.

In the case of domain-specific safety for high-stakes deployment, Omiye et al. \cite{omiye2023racebased} report how LLMs reproduce race-based medical claims proven to be false. Zack et al. \cite{zack2024assessing} audit LLMs for racial and gender stereotypes in clinical recommendations, finding consistently high stereotypical behaviour across all four investigated areas of clinical decision support.



\subsection{Precedents of Demographics-driven Quality-of-Safeguards Disparities}



Beyond general quality-of-service differences between demographic groups, prior work also presents precedents where these differences manifest in the quality of safeguards. On the line of multilingual safety, Yong et al. \cite{yong2023low-resource-ec3} report how translating unsafe prompts into low-resource languages decreases refusal rates by a large margin, revealing an unequally distributed safety coverage across languages. This finding replicates under both unintentional and deliberate attack settings, posing risks to users of low-resource languages as well as through attackers who exploit the jailbreak vulnerability \cite{deng2023multilingual-1fc}. In a follow-up work, Yong et al. \cite{yong-etal-2025-state} quantify this gap through a large-scale survey of the safety literature, finding a widening gap between monolingual and multilingual considerations in safety work between 2020 and 2024 and urging work to address it, as we will do in Chapter \ref{chap:translations}.


On the front of identity-driven disparities, three prior works highlight the need for deeper investigations, of which the first two concern the user-as-victim case and the third targets the user-as-adversary case:

First, Li et al. \cite{li-etal-2024-chatgpt-doesnt} investigate refusal rates for politically sensitive or censored information, conditioned on user biographies disclosed in the prompt. They find younger, female, and Asian-American user personas to evoke higher refusal rates than their counterparts. Refusals are also sycophantic, in the sense that the model is more likely to refuse political positions the user is likely to disagree with.

Second,  Plaza-del-Arco et al. \cite{plaza-del-arco-etal-2025-yes} investigate false-refusal rates based on persona prompting across a large array of model families and sizes. While they cannot establish strong directional patterns across these factors, they find that certain tasks and model families exhibit concerning rates of stereotypical false refusals in the presence of certain gender, ethnicity, and religion personas.

Lastly, Khorramrouz \& Levy \cite{khorramrouz-levy-2026-characterizing} show, on the user-as-adversary front, that models refuse toxic prompts selectively depending on the prompt subject's demographics, and that this selectivity opens an attack vector: a prompt refused for one demographic group can be circumvented by reformulating it around a group for which the model does not refuse.


These identity-based findings motivate our case studies in Chapter \ref{chap:contextualising} investigating identity-based disparities in high-stakes AI-user interactions.









\section{The Science of Evaluations and Construct Validity}


The previous section covered work on \textit{what} is and should be evaluated. Recently, however, the debate has expanded to \textit{how} to evaluate, calling for a \textit{Science of Evaluations} \cite{weidinger2025evaluationscience}.


\subsection{Measurement, Construct Validity, and the Limits of Current Methods}


Already in 2021, Raji et al. \cite{raji2021everything} critiqued AI benchmarks as "construct-invalid stand-ins", claiming that benchmarks cannot be context-independent and hence do not allow general claims about the safety of a model, but only claims about clearly scoped and contextualised sub-categories of the risks. Similarly, Ethayarajh \& Jurafsky \cite{ethayarajh-jurafsky-2020-utility} argue that aggregate leaderboard scores obscure differential fairness outcomes.
Jacobs \& Wallach \cite{jacobs2021measurement} claim that in many cases fairness-related harms stem from a mismatch between the construct intended and what is actually measured. Applied to the case of user-dependent safety, their measurement theoretic approach supports that the underlying claim of making models safe for all often fails to be operationalised in a valid way. Despite these early papers, the challenge of constructing evaluations for generative AI remains four years later: Wallach et al. \cite{wallach2025position} argue that such evaluations are at their core social-science measurements that should follow the process of defining the construct first, then building valid and reliable instruments to measure it. Empirically, these claims are supported by Feffer et al.'s analysis of red-teaming methods, which they find to be often under-specified and performative, lacking shared standards and construct validity \cite{feffer2024redteaming}. Beyond red-teaming, a systematic review of 144 open safety datasets confirms highly idiosyncratic evaluation practices and a clear scarcity of non-English and naturalistic data \cite{rottger2025safetyprompts, grey2025safetymeasurementsystematicliterature}. Similarly, Weidinger et al. \cite{Weidinger2023SociotechnicalSE} critique the skew towards capability evaluations over human-AI-interaction and systemic-impact evaluations, further supporting the need for evaluations that capture human interaction and systemic impact, not only model capability.






\subsection{Reliability and Bias of Human and AI Annotators}


The practice of using LLMs as judges in evaluations is now standard, yet carries documented position, verbosity and self-enhancement biases \cite{NEURIPS2023_91f18a12}, which have been catalogued and surveyed since \cite{gu2025surveyllmasajudge}. LLM-judges often encode distinct, quantifiable biases \cite{ye2025justice}, agree among themselves while diverging from domain experts on expert tasks \cite{limitations_llm_judge}, and exhibit questionable reliability across repeated judgements \cite{schroeder2025trustllmjudgmentsreliability}. The validity of a judgement often hinges on the degrees of freedom researchers have in designing prompts and choosing the judge model and temperature, which can introduce hidden effects into the results \cite{baumann2025largelanguagemodelhacking}.

Not only are LLM-judges potentially biased in their judgement, but also human annotators are. Perceptions of safety and alignment  are culturally and demographically dependent \cite{kirk2024prism}. As shown by Aroyo et al. \cite{aroyo2023dices}, safety ratings vary systematically across the annotators' ethnic, age, and gender distribution. Motivated by this, we will also explore judge-free evaluation settings in Chapter \ref{chap:probes}.




\section{Positioning: Methodological Safety Discrimination}

In the first part of this literature review, we identified that quality-of-service deprivations based on demographics are known and prevalent, but that quality-of-safeguard deprivations can so far be traced only to scattered case studies. In the second part, we discussed the measurement challenge of safety evaluations. This thesis bridges both strands by systematically studying how to evaluate demographics-driven quality-of-safeguards deprivations.

We tackle this in four steps: in Chapter \ref{chap:evaleval}, we will map demographic coverage across 23 safety benchmarks; Chapter \ref{chap:translations} will discuss approaches to multilingual and multidialectal evaluation and faithful translation; Chapter \ref{chap:contextualising} investigates demographic contextualisation of evaluations of advice at the judge, prompt and memory level; and Chapter \ref{chap:probes} explores using cross-lingual linear probes that measure safety-relevant behaviour without relying on LLM-judges.

Together, the case studies and methods presented throughout these chapters establish the ground of methodological safety discrimination and help future work to measure behavioural safety discrimination.
